\begin{ouverture}[courbe]{19}{La proportionnalité en actes}{Partie 5 · Modéliser et organiser}
{Fois plus · Fois moins · Ajouter des situations}
\objectif{Reconnaître une situation de proportionnalité et la justifier.}
\objectif{Raisonner par « fois plus », « fois moins » et par addition de situations.}
\objectif{Résoudre des problèmes en plusieurs étapes.}
\prerequis{Proportionnalité du CM1, multiplication et division.}
\end{ouverture}
\pagestyle{otMain}\essentiel\fontsize{11}{15}\selectfont
\phantomsection\label{cours-19-1}
\essSection{1}{Reconnaître, toujours justifier}
\begin{definition}Deux grandeurs sont en situation de proportionnalité quand multiplier l’une par un nombre multiplie l’autre par le même nombre : 3 fois plus de baguettes, 3 fois plus cher ; 4 fois moins de feuilles, pile 4 fois moins épaisse.\end{definition}
\begin{exemple}Le prix de cahiers identiques est proportionnel à leur nombre. Contre-exemples : la taille selon l’âge, un forfait avec part fixe, un tarif dégressif. La phrase « fois plus » doit être vraie pour tous les nombres, pas seulement une fois.\end{exemple}
\begin{attention}Au cours moyen, le raisonnement s’écrit en français : « Si j’achète 3 fois plus, je paie 3 fois plus. » Ni tableau de proportionnalité, ni produit en croix : les phrases suffisent et montrent le sens.\end{attention}
\begin{savoirfaire}[Le bon sens d’abord]Contrôler la vraisemblance : plus de croissants doit coûter plus cher, moins de temps doit donner moins d’eau. Un résultat qui varie à contresens signale une erreur de raisonnement.\end{savoirfaire}
\clearpage
\phantomsection\label{cours-19-2}
\essSection{2}{Combiner les situations}
\begin{savoirfaire}[Additionner deux situations]Quand 4 stylos coûtent 6 euros et 2 stylos coûtent 3 euros, alors $4+2=6$ stylos coûtent $6+3=9$ euros : deux situations s’additionnent. Avec « fois plus » et « addition », presque tous les nombres deviennent accessibles.\end{savoirfaire}
\begin{exemple}Prix de 10 m de ruban connaissant 4 m à 6 euros : 2 m coûtent 3 euros (2 fois moins), et $10=4+4+2$ : $6+6+3=15$ euros. On enchaîne « fois moins » puis sommes.\end{exemple}
\begin{propriete}Les recettes, les débits réguliers, les prix au mètre ou au kilogramme se traitent ainsi. La donnée de départ sert de brique : on la double, la triple, la partage, puis on assemble.\end{propriete}
\begin{attention}La fraction aide aussi : 6 crêpes utilisent 150 g de farine, donc 1 crêpe environ 25 g ? Oui, mais seulement si la situation est proportionnelle : toujours le vérifier avant de calculer.\end{attention}

\clearpage\rubrique{Découverte}
\begin{blocLivre}{ex-19-1}\begin{exercice}[1]{Proportionnel ou pas ?}
Dire si la situation relève de la proportionnalité : le prix de baguettes identiques ; l’âge et la pointure ; la distance parcourue à allure régulière selon la durée.
\end{exercice}\end{blocLivre}
\begin{blocLivre}{ex-19-2}\begin{exercice}[1]{Fois plus}
Quatre croissants coûtent 5 euros. Combien coûtent 12 croissants ?
\end{exercice}\end{blocLivre}
\begin{blocLivre}{ex-19-3}\begin{exercice}[1]{Fois moins}
Six litres de peinture couvrent 30 m². Quelle surface couvrent 2 litres ? Et 18 litres ?
\end{exercice}\end{blocLivre}
\begin{blocLivre}{ex-19-4}\begin{exercice}[1]{Additionner deux situations}
Quatre stylos coûtent 6 euros. Combien coûtent 2 stylos ? En déduire le prix de 6 stylos.
\end{exercice}\end{blocLivre}
\begin{blocLivre}{ex-19-5}\begin{exercice}[1]{La recette au départ}
Pour 6 crêpes, il faut 150 g de farine. Quelle masse pour 18 crêpes ? Pour 9 crêpes ?
\end{exercice}\end{blocLivre}

\rubrique{Entraînement}
\begin{blocLivre}{ex-19-6}\begin{exercice}[2]{Le débit de la pompe}
Une pompe régulière remplit 80 L en 4 minutes.\begin{questions}
\item Quel volume en 12 minutes ?
\item En combien de temps remplit-elle 200 L ?
\end{questions}
\end{exercice}\end{blocLivre}
\begin{blocLivre}{ex-19-7}\begin{exercice}[2]{Le prix au kilogramme}
Trois kilogrammes de pommes coûtent $5{,}40$ euros. Combien coûte 1 kg ? Et 7 kg ?
\end{exercice}\end{blocLivre}
\begin{blocLivre}{ex-19-8}\begin{exercice}[2]{L’échelle du plan}
Sur un plan, 1 cm représente 50 m réels. Quelle distance réelle pour 7 cm sur le plan ? Et quelle longueur sur le plan pour 350 m réels ?
\end{exercice}\end{blocLivre}
\begin{blocLivre}{ex-19-9}\begin{exercice}[2]{Le piège du forfait}
Un escalade-parc coûte 8 euros d’entrée puis 2 euros par activité. Léo dit : « 6 activités coûtent deux fois plus cher que 3 activités. » Vérifier et conclure.
\end{exercice}\end{blocLivre}
\begin{blocLivre}{ex-19-10}\begin{exercice}[2]{Les deux imprimantes}
L’imprimante A sort 120 pages en 3 minutes ; la B, 150 pages en 4 minutes, toutes deux à allure régulière. Laquelle est la plus rapide ? Justifier en comparant sur 12 minutes.
\end{exercice}\end{blocLivre}

\rubrique{Problèmes}
\begin{blocLivre}{ex-19-11}\begin{exercice}[3]{La course régulière}
Un cycliste roule à allure régulière : 24 km en 1 heure. Quelle distance en 2 h 30 ?
\end{exercice}\end{blocLivre}
\begin{blocLivre}{ex-19-12}\begin{exercice}[3]{Les photocopies de l’école}
Cent photocopies coûtent 4 euros. Combien coûtent 250 photocopies ?
\end{exercice}\end{blocLivre}
\begin{blocLivre}{ex-19-13}\begin{exercice}[3]{La maquette}
Une maquette réduit toutes les longueurs de la même façon : la voiture réelle de 4 m est représentée par 8 cm. Quelle longueur sur la maquette pour un bus de 12 m ? Et pour un vélo de 2 m ?
\end{exercice}\end{blocLivre}
\begin{blocLivre}{ex-19-14}\begin{exercice}[3]{Comparer les lots}
Le lot A offre 5 cahiers pour 7 euros ; le lot B, 10 cahiers pour $13{,}50$ euros. Quel lot est le plus avantageux ?
\end{exercice}\end{blocLivre}
\begin{blocLivre}{ex-19-15}\begin{exercice}[3]{Proportionnel, le périmètre ?}
Quand on triple le côté d’un carré, son périmètre triple-t-il ? Et son aire ? Tester avec un côté de 2 cm, puis conclure.
\end{exercice}\end{blocLivre}

\clearpage
